% GENERATED FILE. Edit canonical lessons, then run npm run generate:books.
% canonical-source-hash: fnv1a64:dcd1d911d4c921d0
% canonical-lessons: JA-C121-shikakui, JA-C121-momoiro, JA-C121-haiiro, JA-C121-usui, JA-C121-asai

\chapter{Square, Pink, Grey, Thin, Shallow}
\label{ch:ja-square-pink-grey-thin-shallow}
\hlchaptermodality{121}

\begin{chapteropening}
\textbf{By the end of this chapter:} \emph{I can say square, pink, grey, thin and shallow.}

Complete the last lesson of chapter 121: i can say square, pink, grey, thin and shallow.
\end{chapteropening}

\section[\texorpdfstring{\ja{しかくい} (shikakui)}{しかくい (shikakui)}]{\ja{しかくい} (shikakui) — square}

\label{lesson:JA-C121-shikakui}



\begin{quote}
Before the new one: say the Japanese for kind, then the Japanese for polite, careful.
\end{quote}

\subsection*{You'll want to know: \ja{しかくい}}
\textbf{\ja{しかくい}} — \emph{shikakui} — \textquotedblleft{}square\textquotedblright{}.

\begin{grammarlens}[title={What you've built}]
One more word for this chapter.
\end{grammarlens}

\subsection*{Your turn}
\begin{itemize}
\raggedright
  \item \emph{Say it:} \emph{shikakui}
  \item \emph{Say it:} \emph{shikakui}, once more
  \item \emph{Say it:} say \emph{shikakui}
  \item \emph{Recall:} say the Japanese for kind, then the Japanese for polite, careful, then say \emph{shikakui} again
\end{itemize}

\begin{tcolorbox}[breakable,colback=teal!4,colframe=teal!35!black,title={Before you move on}]
What does \ja{しかくい} mean? (\textquotedblleft{}Square\textquotedblright{}.) Say it once more.
\end{tcolorbox}
\section[\texorpdfstring{\ja{ももいろ} (momoiro)}{ももいろ (momoiro)}]{\ja{ももいろ} (momoiro) — pink}

\label{lesson:JA-C121-momoiro}



\begin{quote}
Before the new one: say the Japanese for polite, careful, then the Japanese for square.
\end{quote}

\subsection*{You'll want to know: \ja{ももいろ}}
\textbf{\ja{ももいろ}} — \emph{momoiro} — \textquotedblleft{}pink\textquotedblright{}.

\begin{grammarlens}[title={What you've built}]
One more word for this chapter.
\end{grammarlens}

\subsection*{Your turn}
\begin{itemize}
\raggedright
  \item \emph{Say it:} \emph{momoiro}
  \item \emph{Say it:} \emph{momoiro}, once more
  \item \emph{Say it:} say \emph{momoiro}
  \item \emph{Recall:} say the Japanese for polite, careful, then the Japanese for square, then say \emph{momoiro} again
\end{itemize}

\begin{tcolorbox}[breakable,colback=teal!4,colframe=teal!35!black,title={Before you move on}]
What does \ja{ももいろ} mean? (\textquotedblleft{}Pink\textquotedblright{}.) Say it once more.
\end{tcolorbox}
\section[\texorpdfstring{\ja{はいいろ} (haiiro)}{はいいろ (haiiro)}]{\ja{はいいろ} (haiiro) — grey}

\label{lesson:JA-C121-haiiro}



\begin{quote}
Before the new one: say the Japanese for square, then the Japanese for pink.
\end{quote}

\subsection*{You'll want to know: \ja{はいいろ}}
\textbf{\ja{はいいろ}} — \emph{haiiro} — \textquotedblleft{}grey\textquotedblright{}.

\begin{grammarlens}[title={What you've built}]
One more word for this chapter.
\end{grammarlens}

\subsection*{Your turn}
\begin{itemize}
\raggedright
  \item \emph{Say it:} \emph{haiiro}
  \item \emph{Say it:} \emph{haiiro}, once more
  \item \emph{Say it:} say \emph{haiiro}
  \item \emph{Recall:} say the Japanese for square, then the Japanese for pink, then say \emph{haiiro} again
\end{itemize}

\begin{tcolorbox}[breakable,colback=teal!4,colframe=teal!35!black,title={Before you move on}]
What does \ja{はいいろ} mean? (\textquotedblleft{}Grey\textquotedblright{}.) Say it once more.
\end{tcolorbox}
\section[\texorpdfstring{\ja{うすい} (usui)}{うすい (usui)}]{\ja{うすい} (usui) — thin, weak (of taste or colour)}

\label{lesson:JA-C121-usui}



\begin{quote}
Before the new one: say the Japanese for pink, then the Japanese for grey.
\end{quote}

\subsection*{You'll want to know: \ja{うすい}}
\textbf{\ja{うすい}} — \emph{usui} — \textquotedblleft{}thin, weak (of taste or colour)\textquotedblright{}.

\begin{grammarlens}[title={What you've built}]
One more word for this chapter.
\end{grammarlens}

\subsection*{Your turn}
\begin{itemize}
\raggedright
  \item \emph{Say it:} \emph{usui}
  \item \emph{Say it:} \emph{usui}, once more
  \item \emph{Say it:} say \emph{usui}
  \item \emph{Recall:} say the Japanese for pink, then the Japanese for grey, then say \emph{usui} again
\end{itemize}

\begin{tcolorbox}[breakable,colback=teal!4,colframe=teal!35!black,title={Before you move on}]
What does \ja{うすい} mean? (\textquotedblleft{}Thin, weak (of taste or colour)\textquotedblright{}.) Say it once more.
\end{tcolorbox}
\section[\texorpdfstring{\ja{あさい} (asai)}{あさい (asai)}]{\ja{あさい} (asai) — shallow}

\label{lesson:JA-C121-asai}



\begin{quote}
Before the new one: say the Japanese for grey, then the Japanese for thin.
\end{quote}

\subsection*{You'll want to know: \ja{あさい}}
\textbf{\ja{あさい}} — \emph{asai} — \textquotedblleft{}shallow\textquotedblright{}.

\begin{grammarlens}[title={What you've built}]
One more word for this chapter.
\end{grammarlens}

\subsection*{Your turn}
\begin{itemize}
\raggedright
  \item \emph{Say it:} \emph{asai}
  \item \emph{Say it:} \emph{asai}, once more
  \item \emph{Say it:} say \emph{asai}
  \item \emph{Recall:} say the Japanese for grey, then the Japanese for thin, then say \emph{asai} again
\end{itemize}

\begin{tcolorbox}[breakable,colback=teal!4,colframe=teal!35!black,title={Before you move on}]
What does \ja{あさい} mean? (\textquotedblleft{}Shallow\textquotedblright{}.) Say it once more.
\end{tcolorbox}
